\chapter{Leave-One-Subject-Out Fold Results}
\label{app:loso}

% SOURCE: appendix_B_loso_fold_results_v2.md
% STATUS: converted from markdown
% NOTE: [source:] tags stripped per ASSEMBLY_PLAN.md (harvested to appF)


This appendix details the fold-by-fold Leave-One-Subject-Out (LOSO) cross-validation results for the temporal sequence models evaluated in Chapter 13 (\texttt{temporal\_model\_comparison.csv} and \texttt{temporal\_model\_evaluation\_report.md}).



\section{LOSO Protocol Summary}
\label{sec:b-1-loso-protocol-summary}


To evaluate generalization and eliminate subject-identity leakage on small cohorts, Leave-One-Subject-Out (LOSO) cross-validation was executed across:
\begin{itemize}
  \item \textbf{Squat Cohort}: 9 folds corresponding to the 9 subjects ($N=98$ total repetitions).
  \item \textbf{Lunge Cohort}: 7 folds corresponding to the 7 subjects ($N=61$ total repetitions).
\end{itemize}

In each fold, all repetitions of a single held-out subject formed the test set, while the models were trained on the remaining subjects' repetitions.



\section{Fold-by-Fold Results Breakdown}
\label{sec:b-2-fold-by-fold-results-breakdown}


Fold-level LSTM accuracy varies substantially ($20.00\%\text{--}100.00\%$ for squats; $0.00\%\text{--}100.00\%$ for lunges), reflecting instability characteristic of deep sequence models trained on 7–9-subject cohorts. This variance is itself evidence for the endpoint-dominance interpretation reported in Chapter 13: where the peak-flexion baseline maintains consistent per-fold performance, LSTM fold outcomes are dominated by held-out-subject identity rather than form patterns. The globally-pooled aggregates reported in Chapter 13 Table 13.1 weight each fold by test-set rep count and represent the mathematically rigorous summary.

\subsection{Table C.1: Squat Cohort LOSO Fold-by-Fold Results (9 Folds, N=98 Reps)}
\label{sec:table-b-1a-squat-cohort-loso-fold-by-fold-results-9-folds-n-98-reps}


\begin{table}[htbp]
\centering
\scriptsize
\setlength{\tabcolsep}{3pt}
\renewcommand{\arraystretch}{1.2}
\caption{Squat Cohort LOSO Fold-by-Fold Results (9 Folds, N=98 Reps)}
\label{tab:loso-squat-folds}
\begin{tabular}{@{} >{\centering\arraybackslash}p{1.1cm} p{3.0cm} >{\centering\arraybackslash}p{1.8cm} >{\centering\arraybackslash}p{2.0cm} >{\centering\arraybackslash}p{1.8cm} >{\centering\arraybackslash}p{1.8cm} >{\centering\arraybackslash}p{1.8cm}@{}}
\hline
\textbf{Fold} & \textbf{Held-Out Subject} & \textbf{Reps (C/I)} & \textbf{Peak-Flex Baseline} & \textbf{Shape Baseline} & \textbf{LSTM Scheme A} & \textbf{LSTM Scheme B} \\
\hline
Fold 1 & Subject 1 (\texttt{PM\_008}) & 17 (16C / 1I) & 82.35\% & 52.94\% & 94.12\% & 94.12\% \\
Fold 2 & Subject 2 (\texttt{PM\_022}) & 10 (10C / 0I) & 50.00\% & 30.00\% & 60.00\% & 100.00\% \\
Fold 3 & Subject 3 (\texttt{PM\_029}) & 10 (5C / 5I) & 90.00\% & 0.00\% & 60.00\% & 50.00\% \\
Fold 4 & Subject 4 (\texttt{PM\_038}) & 10 (10C / 0I) & 60.00\% & 10.00\% & 20.00\% & 100.00\% \\
Fold 5 & Subject 5 (\texttt{PM\_043}) & 10 (5C / 5I) & 100.00\% & 30.00\% & 70.00\% & 50.00\% \\
Fold 6 & Subject 6 (\texttt{PM\_105}) & 10 (5C / 5I) & 100.00\% & 50.00\% & 30.00\% & 50.00\% \\
Fold 7 & Subject 7 (\texttt{PM\_118}) & 10 (5C / 5I) & 70.00\% & 40.00\% & 60.00\% & 50.00\% \\
Fold 8 & Subject 8 (\texttt{PM\_113}) & 10 (5C / 5I) & 90.00\% & 60.00\% & 40.00\% & 50.00\% \\
Fold 9 & Subject 9 (\texttt{PM\_126}) & 11 (11C / 0I) & 90.91\% & 63.64\% & 72.73\% & 100.00\% \\
\hline
\end{tabular}
\end{table}

\subsection{Table C.2: Lunge Cohort LOSO Fold-by-Fold Results (7 Folds, N=61 Reps)}
\label{sec:table-b-1b-lunge-cohort-loso-fold-by-fold-results-7-folds-n-61-reps}


\begin{table}[htbp]
\centering
\scriptsize
\setlength{\tabcolsep}{3pt}
\renewcommand{\arraystretch}{1.2}
\caption{Lunge Cohort LOSO Fold-by-Fold Results (7 Folds, N=61 Reps)}
\label{tab:loso-lunge-folds}
\begin{tabular}{@{} >{\centering\arraybackslash}p{1.1cm} p{3.0cm} >{\centering\arraybackslash}p{1.8cm} >{\centering\arraybackslash}p{2.0cm} >{\centering\arraybackslash}p{1.8cm} >{\centering\arraybackslash}p{1.8cm} >{\centering\arraybackslash}p{1.8cm}@{}}
\hline
\textbf{Fold} & \textbf{Held-Out Subject} & \textbf{Reps (C/I)} & \textbf{Peak-Flex Baseline} & \textbf{Shape Baseline} & \textbf{LSTM Scheme A} & \textbf{LSTM Scheme B} \\
\hline
Fold 1 & Subject 2 (\texttt{PM\_028}) & 9 (4C / 5I) & 77.78\% & 66.67\% & 44.44\% & 55.56\% \\
Fold 2 & Subject 3 (\texttt{PM\_037}) & 10 (0C / 10I) & 40.00\% & 60.00\% & 40.00\% & 0.00\% \\
Fold 3 & Subject 4 (\texttt{PM\_117b}) & 9 (5C / 4I) & 100.00\% & 44.44\% & 55.56\% & 44.44\% \\
Fold 4 & Subject 5 (\texttt{PM\_042}) & 1 (0C / 1I) & 0.00\% & 0.00\% & 100.00\% & 100.00\% \\
Fold 5 & Subject 6 (\texttt{PM\_104}) & 10 (5C / 5I) & 90.00\% & 60.00\% & 50.00\% & 50.00\% \\
Fold 6 & Subject 7 (\texttt{PM\_117a}) & 10 (5C / 5I) & 90.00\% & 50.00\% & 60.00\% & 50.00\% \\
Fold 7 & Subject 9 (\texttt{PM\_125}) & 12 (6C / 6I) & 91.67\% & 66.67\% & 66.67\% & 50.00\% \\
\hline
\end{tabular}
\end{table}



\section{Observations}
\label{sec:b-3-observations}

The fold-by-fold breakdown demonstrates consistency across held-out subjects:
\begin{enumerate}
  \item \textbf{Peak Flexion Stability}: The single-feature peak flexion baseline maintained high classification performance ($50.00\%\text{--}100.00\%$ across individual fold test sets for squats; $40.00\%\text{--}100.00\%$ for lunges) across held-out subjects.
  \item \textbf{LSTM Overfitting Ceiling}: The deep sequence network failed to achieve stable generalization across held-out subject splits, exhibiting high variance under Scheme A and collapsing to majority-class prediction under Scheme B.
\end{enumerate}